% !TeX root = ../navier-stokes-manuscript.tex
\subsection{More spatial regularity, not less}\label{sub:ThePressureCompleteMixedDerivativeTower}

Read only in time, every new rung looks worse.
A nontrivial change squeezed into a shorter interval produces a larger response, and the temporal Tower appears to run straight toward singularity.
The same rows carry a different count in space.
Viscosity inserts a Laplacian every time the equation is used, and at critical height the temporal and spatial scale weights line up.
All identities below are taken on an interval where the solution is classical, so the differentiations and pairings are justified.

For
\[
  E_s(t):=\frac12\norm{\Lambda^su(t)}_2^2,
  \qquad
  E_{s,\lambda}(t):=\frac12\norm{\Lambda^su_\lambda(t)}_2^2,
\]
the calculation from critical height gives
\begin{equation}\label{eq:BodySpatialEnergyScaling}
  E_{s,\lambda}(t)
  =\lambda^{2s-1}E_s(\lambda^2t).
\end{equation}
Since \(H=E_{1/2}\),
\[
  H_\lambda(t)=H(\lambda^2t).
\]
After \(n\) time derivatives,
\begin{equation}\label{eq:BodyCriticalHeightTemporalScaling}
  H_\lambda^{(n)}(t)
  =\lambda^{2n}H^{(n)}(\lambda^2t).
\end{equation}
The spatial energy at height \(n+\tfrac12\) has the same scale factor:
\begin{equation}\label{eq:BodySpatialHeightScalingDiagonal}
  E_{n+1/2,\lambda}(t)
  =\lambda^{2n}E_{n+1/2}(\lambda^2t).
\end{equation}
Equal scale weight does not identify \(H^{(n)}\) with \(E_{n+1/2}\).
It selects the spatial height that belongs beside the \(n\)-th temporal response.

Applying \(\Lambda^s\) to Navier--Stokes and pairing the result with \(\Lambda^su\) gives the viscous contribution
\[
  \nu\pair{\Lambda^su}{\Lambda^s\Delta u}_{L^2}
  =-\nu\norm{\Lambda^{s+1}u}_2^2
  =-2\nu E_{s+1}.
\]
Keep everything outside the viscous term in the complete transport--pressure transfer
\[
  \mathcal P_s
  :=-
  \operatorname{Re}
  \pair{\Lambda^su}
  {\Lambda^s\bigl((u\cdot\nabla)u+\nabla p\bigr)}_{L^2}.
\]
The global pressure pairing vanishes against the divergence-free velocity.
It is retained inside \(\mathcal P_s\) so that the row remains pressure-complete before that cancellation is used.
The exact spatial-energy row is
\begin{equation}\label{eq:BodyCriticalEnergyBalance}
  E_s'=\mathcal P_s-2\nu E_{s+1}.
\end{equation}

At critical height, the first time derivative already reaches the next spatial energy:
\[
  H'=\mathcal P_{1/2}-2\nu E_{3/2}.
\]
Differentiate again and use the same row at \(s=\tfrac32\):
\[
\begin{aligned}
  H''
  &=\partial_t\mathcal P_{1/2}-2\nu E_{3/2}'
  \\
  &=\partial_t\mathcal P_{1/2}
  -2\nu\mathcal P_{3/2}
  +(2\nu)^2E_{5/2}.
\end{aligned}
\]
The third derivative reaches one height farther:
\[
  H'''
  =\partial_t^2\mathcal P_{1/2}
  -2\nu\partial_t\mathcal P_{3/2}
  +(2\nu)^2\mathcal P_{5/2}
  -(2\nu)^3E_{7/2}.
\]
Each differentiation keeps the transfer terms already produced and adds one higher viscous energy.
For \(n\geq1\), induction gives
\begin{equation}\label{eq:BodyCriticalHeightSpatialAscent}
  H^{(n)}
  =(-2\nu)^nE_{n+1/2}
  +\sum_{j=0}^{n-1}
  (-2\nu)^j
  \partial_t^{\,n-1-j}\mathcal P_{j+1/2}.
\end{equation}
Because \(\nu>0\), the top spatial energy can be isolated:
\begin{equation}\label{eq:BodyCompletedCriticalRow}
  E_{n+1/2}
  =(-2\nu)^{-n}
  \left(
    H^{(n)}
    -\sum_{j=0}^{n-1}
      (-2\nu)^j
      \partial_t^{\,n-1-j}\mathcal P_{j+1/2}
  \right).
\end{equation}
The identity exposes the higher Sobolev energy and isolates it algebraically.
It gives no bound until \(H^{(n)}\) and every transfer term beside it have also been controlled.

The temporal Tower enters through \(H^{(n)}\) itself.
Since
\[
  H(t)
  =\frac12
  \pair{\Lambda^{1/2}V_0(t)}
       {\Lambda^{1/2}V_0(t)}_{L^2},
\]
Leibniz's rule gives
\begin{equation}\label{eq:BodyHeightDerivativeTemporalContraction}
  H^{(n)}(t)
  =\frac12\sum_{a=0}^{n}\binom{n}{a}
  \operatorname{Re}
  \pair{\Lambda^{1/2}V_a(t)}
       {\Lambda^{1/2}V_{n-a}(t)}_{L^2}.
\end{equation}
The three entries on one rung remain different objects.
The field response is \(V_n\), the critical response \(H^{(n)}\) contracts the temporal rows, and viscosity exposes the spatial energy \(E_{n+1/2}\).

\Needspace{10\baselineskip}
\begin{center}
\captionsetup{hypcap=false}
\captionof{table}{The first time--space rungs of the critical-height calculation.}
\label{tab:BodyTemporalSpatialDiscovery}
\small
\renewcommand{\arraystretch}{1.2}
\begin{tabular}{@{}c c c c@{}}
\toprule
order & Eulerian response & critical response & spatial height \\
\midrule
\(0\) & \(V_0=u\) & \(H=E_{1/2}\) & \(E_{1/2}\) \\
\(1\) & \(V_1=u_t\) & \(H'\) & \(E_{3/2}\) \\
\(2\) & \(V_2=u_{tt}\) & \(H''\) & \(E_{5/2}\) \\
\(3\) & \(V_3=u_{ttt}\) & \(H'''\) & \(E_{7/2}\) \\
\(n\) & \(V_n=\partial_t^nu\) & \(H^{(n)}\) & \(E_{n+1/2}\) \\
\bottomrule
\end{tabular}
\end{center}

Whenever the Zeno route keeps a nontrivial change alive, the temporal column asks for a larger response.
The matched spatial column exposes higher orders in the Sobolev scale through
\(E_{1/2},E_{3/2},E_{5/2},E_{7/2},\ldots\).

This is still an exposed ladder, not a regularity estimate.
The transfer terms remain unpaid, and one diagonal inside the velocity row does not fill the whole time--space Tower.

Every temporal response needs its own spatial ladder.
For \(r\in\Nzero\), let \(E_{r,s}:=\tfrac12\norm{\Lambda^sV_r}_2^2\), and define the complete transport--pressure transfer in the \(r\)-th recurrence by
\[
  \mathcal P_{r,s}
  :=-
  \operatorname{Re}
  \pair{\Lambda^sV_r}
  {\Lambda^s\!\left[
    \sum_{a=0}^{r}\binom{r}{a}(V_a\cdot\nabla)V_{r-a}
    +\nabla Q_r
  \right]}_{L^2}.
\]
Pairing the \(r\)-th temporal equation with \(\Lambda^{2s}V_r\) gives
\begin{equation}\label{eq:BodyEveryTemporalRowSpatialEnergy}
  E_{r,s}'=\mathcal P_{r,s}-2\nu E_{r,s+1}.
\end{equation}
After any finite ascent depth \(k\geq1\),
\begin{equation}\label{eq:BodyEveryTemporalRowCompletedHeight}
  \partial_t^kE_{r,1/2}
  =(-2\nu)^kE_{r,k+1/2}
  +\sum_{j=0}^{k-1}
  (-2\nu)^j
  \partial_t^{\,k-1-j}\mathcal P_{r,j+1/2}.
\end{equation}
The temporal index \(r\) chooses the response, while the spatial index \(k\) chooses the height inside that response.
They are independent coordinates of the same mixed jet.

Together they form the grid
{\small
\[
\begin{array}{c|ccccc}
 & H^0_x & H^1_x & H^2_x & H^3_x & \cdots\\
\hline
V_0=u      & \bullet & \bullet & \bullet & \bullet & \cdots\\
V_1=u_t    & \bullet & \bullet & \bullet & \bullet & \cdots\\
V_2=u_{tt} & \bullet & \bullet & \bullet & \bullet & \cdots\\
V_3=u_{ttt}& \bullet & \bullet & \bullet & \bullet & \cdots\\
\vdots     & \vdots  & \vdots  & \vdots  & \vdots  & \ddots
\end{array}
\]
}
The bullets mark formal coordinates of the mixed jet, not bounds already proved for them.
Moving down selects a later temporal response, while moving across selects a higher Sobolev height inside that response.
For the fixed classical interval \(I=(0,T)\), controlling every coordinate would place the solution in the mixed intersection
\begin{equation}\label{eq:BodyDesiredMixedIntersection}
  u\in
  \bigcap_{r=0}^{\infty}
  \bigcap_{m=0}^{\infty}
  H^r(I;H^m(\R^3)).
\end{equation}

No finite estimate can cover this infinite grid at once.
For fixed \(N\) and \(M\), the rows through \(N\) and columns through \(M\) form one finite rectangle, and every coordinate belongs to some such rectangle.
The question has now changed form.
Instead of controlling the infinite intersection directly, the next section has to control every finite rectangle in a way that agrees on all overlaps.
